\documentclass{article}
\usepackage{tekvideoexercises}

\begin{document}
\exercisename{De Moivres formel}
\streaklength{5}

\tableofcontents
\newpage

\begin{exercise}{De Moivres formel 1-1}

Skriv $(2(\cos(\frac{\pi}{4}) + i\sin(\frac{\pi}{4})))^3$ i polær form ved brug af De Moivres formel.

\answerbox{8\left(\cos\left(\frac{3\pi}{4}\right) + i\sin\left(\frac{3\pi}{4}\right)\right)}

\hint

De Moivres formel: $(r(\cos \theta + i \sin \theta))^n = r^n(\cos(n\theta) + i\sin(n\theta))$

\hint

Her er $r = 2$, $\theta = \frac{\pi}{4}$, $n = 3$:

\hint

Beregn $r^n = 2^3 = 8$

\hint

Beregn $n\theta = 3\cdot\frac{\pi}{4} = \frac{3\pi}{4}$

\end{exercise}

\newpage

\begin{exercise}{De Moivres formel 1-2}

Skriv $(3(\cos(\frac{\pi}{3}) + i\sin(\frac{\pi}{3})))^2$ i polær form.

\answerbox{9\left(\cos\left(\frac{2\pi}{3}\right) + i\sin\left(\frac{2\pi}{3}\right)\right)}

\hint

Anvend De Moivres formel:

\hint

$r^n = 3^2 = 9$

\hint

$n\theta = 2\cdot\frac{\pi}{3} = \frac{2\pi}{3}$

\end{exercise}

\newpage

\begin{exercise}{De Moivres formel 1-3}

Skriv $(1(\cos(\frac{2\pi}{5}) + i\sin(\frac{2\pi}{5})))^5$ i polær form.

\answerbox{1\left(\cos(2\pi) + i\sin(2\pi)\right)}

\hint

Anvend De Moivres formel:

\hint

$r^n = 1^5 = 1$

\hint

$n\theta = 5\cdot\frac{2\pi}{5} = 2\pi$

\end{exercise}

\newpage

\begin{exercise}{De Moivres formel 1-4}

Skriv $(4(\cos(\frac{\pi}{6}) + i\sin(\frac{\pi}{6})))^4$ i polær form.

\answerbox{256\left(\cos\left(\frac{2\pi}{3}\right) + i\sin\left(\frac{2\pi}{3}\right)\right)}

\hint

Anvend De Moivres formel:

\hint

$r^n = 4^4 = 256$

\hint

$n\theta = 4\cdot\frac{\pi}{6} = \frac{4\pi}{6} = \frac{2\pi}{3}$

\end{exercise}

\newpage

\begin{exercise}{De Moivres formel 1-5}

Skriv $(5(\cos(\frac{3\pi}{4}) + i\sin(\frac{3\pi}{4})))^2$ i polær form.

\answerbox{25\left(\cos\left(\frac{3\pi}{2}\right) + i\sin\left(\frac{3\pi}{2}\right)\right)}

\hint

Anvend De Moivres formel:

\hint

$r^n = 5^2 = 25$

\hint

$n\theta = 2\cdot\frac{3\pi}{4} = \frac{6\pi}{4} = \frac{3\pi}{2}$

\end{exercise}

\newpage

\begin{exercise}{De Moivres formel 1-6}

Skriv $(2(\cos(\frac{\pi}{12}) + i\sin(\frac{\pi}{12})))^6$ i polær form.

\answerbox{64\left(\cos\left(\frac{\pi}{2}\right) + i\sin\left(\frac{\pi}{2}\right)\right)}

\hint

Anvend De Moivres formel:

\hint

$r^n = 2^6 = 64$

\hint

$n\theta = 6\cdot\frac{\pi}{12} = \frac{6\pi}{12} = \frac{\pi}{2}$

\end{exercise}

\end{document}
